\documentclass[11pt,a4paper]{article}
\usepackage{fontspec,amsmath,amssymb,amsthm,microtype}
\setmainfont{lmroman10-regular.otf}[BoldFont=lmroman10-bold.otf,ItalicFont=lmroman10-italic.otf,BoldItalicFont=lmroman10-bolditalic.otf]
\usepackage[margin=25mm]{geometry}
\usepackage[unicode,hidelinks]{hyperref}
\hypersetup{pdfauthor={LangevinLean Project}}
\setlength{\emergencystretch}{3em}
\setlength{\parskip}{0.4em}
\setlength{\parindent}{0pt}
\allowdisplaybreaks[2]
\raggedbottom
\title{\texorpdfstring{Expected energy comparison for the perturbation force}{Expected energy comparison for the perturbation force}}
\author{LangevinLean Project}
\date{}
\begin{document}\maketitle
Draft mathematical exposition; these notes may contain mistakes.


Let \(e\) be convex and differentiable, assume \(e\le f\le e+\delta\), and let \(Y_t\) have the stated stationary smoothed law, integral equation, and future-noise orthogonality. Assume the value Lipschitz bounds needed for integrability. Convex support bounds the force pairing by an envelope-energy difference. Stationarity identifies the time-dependent reference energy with the terminal energy, and the uniform gap then compares this with the original nonsmooth energy in the proximal estimate.


\section{\texorpdfstring{The expected uniform value-\allowbreak{}gap inequality}{The expected uniform value-gap inequality}}


Suppose \(f(x)\le e(x)+\delta\) and both functions are Lipschitz. For \(L^2\) variable \(X\) their values are integrable, so integral monotonicity gives \(\mathbb E f(X)-\mathbb E e(X)\le\delta\). In the notation \(f\le e+\delta\) and an \(L^2\) state, the mass-one calculation is \(\mathbb Ef(X)\le\mathbb Ee(X)+\delta P(\Omega)=\mathbb Ee(X)+\delta\).



The subsequent comparison uses \(e\le f\) at the algorithm’s intermediate state and \(f\le e+\delta\) at the reference endpoint. The integrated gap is \(\delta\) because the source is a probability. All expected values are finite before subtraction, and no derivative of the original nonsmooth function is used.



\section{\texorpdfstring{The stationary convex-\allowbreak{}support and future-\allowbreak{}drift estimates}{The stationary convex-support and future-drift estimates}}


Convex support bounds \(\langle S-Y_h,\nabla e(Y_t)\rangle\) by \(\mathbb E e(S)-\mathbb E e(Y_t)+\langle Y_t-Y_h,\nabla e(Y_t)\rangle\). Stationarity gives equality of \(\mathbb E e(Y_t)\) and \(\mathbb E e(Y_h)\). The future-noise estimate bounds the remaining drift pairing by \(hGK\). The gap \(e\le f\le e+\delta\) changes the difference of the two energies by at most \(\delta\). Integration therefore gives

\[2\left\langle S-Y_h,\int_0^h\nabla e(Y_t)dt\right\rangle-2h(\mathbb E f(S)-\mathbb E f(Y_h))\le2h\delta+2h^2GK.\]

More explicitly, for \(b_t=\nabla e(Y_t)\) in \(L^2\) the support and future-increment bounds give

\[\langle S-Y_h,b_t\rangle\le\mathbb Ee(S)-\mathbb Ee(Y_t)+(h-t)GK.\]

Stationarity replaces \(\mathbb Ee(Y_t)\) by \(\mathbb Ee(Y_h)\). Integration over \([0,h]\) is legitimate because \(b_t\) is continuous and the fixed inner-product functional commutes with the Bochner integral. A valid upper bound is \(h[\mathbb Ee(S)-\mathbb Ee(Y_h)]+h^2GK\).



The decomposition \(S-Y_h=(S-Y_t)+(Y_t-Y_h)\) separates the convex support term from the future-increment term. Again \(\mathbb E e(Y_t)\) equals \(\mathbb E e(Y_h)\). At \(S\) use \(e\le f\), and at \(Y_h\) use \(f\le e+\delta\). Only one additive gap \(\delta\) remains. This is the energy comparison used to cancel the original potential difference in the splitting estimate.





Earlier results used here: \cite[Section~3.13.2]{LangevinLean2026ULASC}; \cite[Section~3.13.5]{LangevinLean2026ULASC}.

\begin{thebibliography}{ULA-LSI}
\bibitem[ULA-SC]{LangevinLean2026ULASC}
LangevinLean Project. \emph{Unadjusted Langevin under Strong Convexity}. LL-01, 2026. Draft report, content revision 7e430ee07530. \url{https://odakyildiz.github.io/LangevinLean/docs/pdf/UnadjustedLangevin/FullProof.pdf}.
\end{thebibliography}

\end{document}
